\section{Productivity Boundary and Research Scope}\label{sec:productivity}
The governing research objective is to reduce the cost and time
required for agents to complete large software projects correctly,
without increasing coordination failures or weakening acceptance
contracts. That objective outruns what any single paper can measure,
and this manuscript is explicit about the ladder of evidence it
occupies.

\paragraph{What this paper is.}
It is a mechanism study of context retirement
(Section~\ref{sec:retire}) plus the verification and evidence
components the surviving rule depends on
(Sections~\ref{sec:model}--\ref{sec:eval}). Two protocol generations
were run; the current one is factorial, pre-registered, and
git-frozen before data generation. Five closure rules were declared
before data; two of the paper's own proposed extensions failed their
gates and are reported as rejected: the byte-regime productivity
claim for frontier retirement (it saves \RetwoBytesDelta{} on
submitted bytes versus masking---no win claimed) and the demand-driven
adaptive-coordination variant ($\RetwoAdaptiveReduction{}$ regime-O
versus plain frontier; promotion gate $\ge$10\%). What survives is
(i)~a necessary conversion rule---retire history only after
converting it into verified bindings, typed obligations, and declared
dependencies; lossy conversion fails the acceptance contract at every
scale tested (28/84 at the largest grid); (ii)~an operations result
at equal correctness---frontier retirement cuts on-demand retrieval
operations \RetwoOpsReduction{} and round-trip-regime cost
\RetwoCostOReduction{} versus observation masking, persisting at
$4\times$ horizon, with a clean capsule audit; and (iii)~a
density-selected conditional role for certificate-tree evidence
layouts (\RetwoDensityRegret{} decision-rule regret). HACT is not the
paper's center; it survives only where its ablation says it wins.

\paragraph{The boundary that remains.}
The retirement study is mechanism-level. It models information flow
between a context policy and a seeded work stream; it does not call a
hosted model, and it cannot establish that any agent finishes faster,
bills less, or reasons better. Its staleness, obligation-loss, and
audit results are structural predictions about policies, not
measurements of agents. The live-agent cohort that could test those
predictions has its protocol pre-registered in
Section~\ref{sec:live-protocol} and has not been run. The negative
controls below---flat checker timings, per-episode verification
overhead, unchanged root-status inputs---are reported as limits of
the serialization claim, and the retirement cohort's own closures are
reported as limits of the mechanism claim. Nothing in this paper is
an agent-productivity measurement, and the abstract says so.
